\section{Parcial 2018 Q1}

% =====================================================================
\subsection{Ejercicio 1 — Secreto compartido de Shamir}
% =====================================================================

\begin{consigna}
Considerar un esquema de secreto compartido de Shamir $(3,3)$ en $\mathbb{Z}_{11}$, con el
conjunto de sombras $C = \{(1,2),(2,6),(3,5)\}$.
\begin{enumerate}[label=\alph*)]
  \item Recuperar el secreto.
  \item Se desea, manteniendo las sombras existentes, repartir más sombras que permitan
        recuperar el secreto.
        \begin{enumerate}[label=\roman*.]
          \item Generar una.
          \item ¿Cuántas distintas se pueden generar?
        \end{enumerate}
\end{enumerate}
\end{consigna}

\paragraph{Temas.} Clase 8 — Principios de Diseño y Vulnerabilidades (\S Secreto compartido de
Shamir en términos de entropía); base de \emph{flujo de información} (Clase 9 — Flujo de
Información y Malware, \S Entropía).

\paragraph{Qué necesitás saber.} Un esquema $(T,N)$ de Shamir reparte el secreto $S$ entre
$N$ sombras de modo que se necesitan \textbf{al menos $T$} para reconstruirlo, y con menos de
$T$ la entropía del secreto \emph{no cambia} (no hay flujo de información). El esquema se
construye con un polinomio de grado $T-1$ sobre un cuerpo finito $\mathbb{Z}_p$:
\begin{equation*}
  f(x) = a_0 + a_1 x + a_2 x^2 + \dots + a_{T-1} x^{T-1} \pmod{p},
\end{equation*}
donde el secreto es el término independiente $S = a_0 = f(0)$ y cada sombra es un par
$(x_i,\, y_i=f(x_i))$. Con $T$ puntos, el polinomio queda unívocamente determinado (un
polinomio de grado $T-1$ se fija con $T$ puntos), y se recupera $S=f(0)$ por
\textbf{interpolación de Lagrange}:
\begin{equation*}
  f(0) = \sum_{j} y_j \prod_{m \neq j} \frac{0 - x_m}{x_j - x_m} \pmod{p}.
\end{equation*}

\paragraph{Notas y conexiones.} La validez del esquema se demuestra con entropía
(Clase 8): $H(S\mid S_1)=H(S)$, $H(S\mid S_1,S_2)=H(S)$ (con $<T$ sombras no hay flujo) y
$H(S\mid S_1,S_2,S_3)=0$ (con $T$ sombras la incertidumbre cae a $0$ y el secreto queda
determinado). Aquí $T=N=3$: hacen falta las tres sombras, ninguna sombra propia revela nada del
secreto. Toda la aritmética es \textbf{módulo $p=11$} (cuerpo finito), incluidos los inversos
multiplicativos para dividir.

\paragraph{Resolución.}
\textbf{a) Recuperar el secreto.} Polinomio de grado $T-1=2$. Interpolación de Lagrange en
$x=0$ con los puntos $(1,2),(2,6),(3,5)$ sobre $\mathbb{Z}_{11}$. Los coeficientes de Lagrange
en $0$ son:
\begin{align*}
  \ell_1(0) &= \frac{(0-2)(0-3)}{(1-2)(1-3)} = \frac{6}{2} = 3, \\
  \ell_2(0) &= \frac{(0-1)(0-3)}{(2-1)(2-3)} = \frac{3}{-1} = -3 \equiv 8, \\
  \ell_3(0) &= \frac{(0-1)(0-2)}{(3-1)(3-2)} = \frac{2}{2} = 1.
\end{align*}
Entonces
\begin{equation*}
  S = f(0) = 2\cdot 3 + 6\cdot 8 + 5\cdot 1 = 6 + 48 + 5 = 59 \equiv 4 \pmod{11}.
\end{equation*}
\fbox{$S = 4$.} (Verificación reconstruyendo el polinomio completo: $f(x)=4+6x+3x^2 \bmod 11$,
que en efecto da $f(1)=13\equiv2$, $f(2)=28\equiv6$, $f(3)=49\equiv5$.)

\medskip
\textbf{b.i) Generar una sombra nueva.} Manteniendo las sombras existentes hay que usar
\textbf{el mismo polinomio} $f(x)=4+6x+3x^2 \bmod 11$ y evaluarlo en un $x$ no usado todavía
(y distinto de $0$, que es el secreto mismo). Por ejemplo, en $x=4$:
\begin{equation*}
  f(4) = 4 + 6\cdot4 + 3\cdot16 = 4 + 24 + 48 = 76 \equiv 10 \pmod{11}.
\end{equation*}
Una sombra nueva válida es $\boxed{(4,\,10)}$. Cualquier subconjunto de $3$ sombras (viejas o
nuevas) reconstruye $S=4$.

\medskip
\textbf{b.ii) ¿Cuántas distintas se pueden generar?} Las sombras se identifican por su
abscisa $x \in \mathbb{Z}_{11}$. Quedan excluidas $x=0$ (es $f(0)=S$, no es una sombra: la
revelaría) y las tres ya repartidas $x\in\{1,2,3\}$. Restan
\begin{equation*}
  11 - 1 - 3 = \boxed{7} \quad\text{sombras nuevas distintas } (x \in \{4,5,6,7,8,9,10\}),
\end{equation*}
con valores $f(4)=10,\ f(5)=10,\ f(6)=5,\ f(7)=6,\ f(8)=2,\ f(9)=4,\ f(10)=1$. El cuerpo
$\mathbb{Z}_{11}$ acota el total de identificadores posibles a $10$ (sin contar el $0$), de los
cuales $3$ ya están usados.

% =====================================================================
\subsection{Ejercicio 2 — Verdadero o falso (control de acceso, OIDC, inyección)}
% =====================================================================

\begin{consigna}
Responder verdadero o falso, justificando o corrigiendo como fuere apropiado:
\begin{enumerate}[label=\alph*)]
  \item Para permitir que usuarios anónimos accedan a un sistema, en vez de utilizar una lista
        de capacidades CAP, es más práctico implementar una lista de control de accesos ACL.
  \item OpenID Connect permite que mediante un password se pueda autorizar a un usuario a
        loguearse en el sistema.
  \item Una de las defensas más recomendadas para evitar ataques de SQL injection es escapar
        los datos ingresados por el usuario (como las comillas).
  \item XSS es un aprovechamiento malicioso de la confianza que un sitio web tiene en la sesión
        que estableció con el browser del usuario.
\end{enumerate}
\end{consigna}

\paragraph{Temas.} Clase 6 — Políticas y Control de Acceso (\S ACL vs.\ capability lists);
Clase 7 — Autenticación (\S Autenticación $\neq$ Autorización, OAuth2/OIDC/SAML);
Clase 8 — Principios de Diseño y Vulnerabilidades (\S Catálogo de vulnerabilidades:
CWE-89 SQLi, CWE-79 XSS; CSRF).

\paragraph{Qué necesitás saber.}
\begin{itemize}[leftmargin=1.4em,itemsep=2pt]
  \item \textbf{ACL vs.\ capability:} la matriz de accesos tiene sujetos en las filas y objetos
        en las columnas. La \textbf{ACL} es una \emph{columna} (por objeto: lista de
        $(sujeto, permisos)$); la \textbf{capability list} es una \emph{fila} (por sujeto: qué
        puede hacer ese sujeto en cada objeto).
  \item \textbf{OAuth2 = autorización; OIDC = autenticación} (capa de identidad sobre OAuth2,
        \emph{quién} sos); \textbf{SAML} = intercambio IdP$\leftrightarrow$SP. Autenticación
        ($\neq$) Autorización.
  \item \textbf{SQLi (CWE-89):} input no controlado que ejecuta SQL; mitigación canónica
        \emph{prepared statements / parametrizar}; escapar es paliativo. \textbf{XSS (CWE-79):}
        inyectar JS vía un input que se renderiza como HTML y lo ejecuta \emph{la víctima}.
\end{itemize}

\paragraph{Notas y conexiones.} El enunciado d) describe en realidad \textbf{CSRF}
(\emph{Cross-Site Request Forgery}), no XSS: CSRF abusa de la confianza del \emph{sitio} en la
sesión/cookies del usuario (el sitio confía en el browser autenticado); XSS abusa de la
confianza del \emph{usuario} en el sitio (el browser ejecuta script malicioso servido por el
sitio). Es la confusión clásica que buscan en el parcial. Inyección (mezclar datos con código)
es $\sim$90\% de las vulnerabilidades (Clase 8).

\paragraph{Resolución.}
\textbf{a) FALSO.} Está al revés. La ACL lista, \emph{por objeto}, los pares
$(sujeto,permisos)$: requiere \textbf{identificar} a cada sujeto, justo lo que no se tiene con
usuarios anónimos. Para acceso anónimo conviene la \textbf{capability}: el usuario presenta una
credencial/token (la capacidad) que \emph{en sí misma} habilita el acceso, sin que el sistema
deba conocer su identidad ni tenerlo enumerado en la lista del objeto. (Además, las capabilities
son el modelo dual, más cómodo en sistemas distribuidos.)

\medskip
\textbf{b) FALSO.} Mezcla dos cosas. \textbf{OpenID Connect autentica}, no autoriza: es la capa
de \emph{autenticación} sobre OAuth2 (saber \emph{quién} es el usuario, ``login con Google''),
y justamente su gracia es \textbf{no} manejar el password en el sistema relyente, sino delegarlo
al IdP. El que \emph{autoriza} (delega acceso a recursos) es \textbf{OAuth2}. Corrección:
``OpenID Connect permite \emph{autenticar} a un usuario delegando el login en un IdP, sin que el
sistema maneje su password''.

\medskip
\textbf{c) VERDADERO} (con matiz). Escapar/sanitizar la entrada (las comillas, etc.) evita que
el input rompa la sintaxis SQL, así que es una defensa recomendada. \textbf{Pero la defensa
canónica y más robusta son los \emph{prepared statements} (parametrizar)}, que separan
estructuralmente el código SQL de los datos y no dependen de acertar todos los caracteres a
escapar. Escapar correctamente es difícil y propenso a errores; parametrizar lo resuelve de
raíz.

\medskip
\textbf{d) FALSO.} Eso describe \textbf{CSRF}, no XSS. CSRF explota la confianza del
\emph{servidor} en la sesión ya autenticada del usuario (sus cookies viajan automáticamente).
\textbf{XSS} (CWE-79) es lo inverso: el atacante inyecta script en un input que el sitio
renderiza como HTML, y ese script se ejecuta en el browser de \emph{la víctima} explotando la
confianza del \emph{usuario} en el sitio. Corrección: ``XSS es la inyección de código (JS) que
el sitio renderiza y ejecuta en el browser de la víctima; se mitiga escapando/sanitizando toda
la salida''.

% =====================================================================
\subsection{Ejercicio 3 — Política de integridad sobre un retículo (Biba)}
% =====================================================================

\begin{consigna}
Un sistema simple implementa una política de seguridad que hace uso del siguiente retículo:
\begin{center}
\begin{tikzpicture}[font=\small,>=Stealth,node distance=10mm]
  \node (ceo) at (0,2) {CEO};
  \node (g1)  at (-1.8,0.6) {Gerente1};
  \node (g2)  at (1.8,0.6)  {Gerente2};
  \node (emp) at (0,-0.8) {Empleado};
  \draw[->] (g1) -- (ceo);
  \draw[->] (g2) -- (ceo);
  \draw[->] (emp) -- (g1);
  \draw[->] (emp) -- (g2);
\end{tikzpicture}
\end{center}
\begin{enumerate}[label=\alph*)]
  \item Sabiendo que $i(\text{Gerente1}) \le i(\text{CEO})$, etc.\ (es decir, las flechas
        indican el flujo de información en términos de una relación de dominancia), indicar:
        \begin{enumerate}[label=\roman*.]
          \item De qué política se trata.
          \item Cuáles y en qué consisten (definir) las principales propiedades que busca hacer
                cumplir esta política. Teniendo éstas en cuenta, ¿qué pueden hacer los usuarios
                en base al modo de acceso \texttt{w} (\texttt{write}, escritura)?
        \end{enumerate}
  \item Si se invierte la relación (el sentido de las flechas), ¿qué política se estaría
        modelando? Compararla brevemente con la política del punto anterior.
\end{enumerate}
\end{consigna}

\paragraph{Temas.} Clase 6 — Políticas y Control de Acceso (\S Modelo Biba (integridad);
\S Compartimentos, reticulado y dominancia; Bell--LaPadula).

\paragraph{Qué necesitás saber.} La etiqueta usada es $i(\cdot)$, un \textbf{nivel de
integridad}: a mayor nivel, mayor confianza. Eso es \textbf{Biba} (integridad), el inverso de
Bell--LaPadula (confidencialidad). En un \textbf{retículo} (orden parcial), una flecha
$A \to B$ con $i(A)\le i(B)$ indica que $B$ \textbf{domina} a $A$ y que la información fluye
``hacia arriba'' siguiendo la flecha. Reglas de \textbf{Biba estricto}:
\begin{align*}
  \textbf{Escritura (Write Down):}\quad & s \text{ puede modificar } o \iff i(s) \ge i(o), \\
  \textbf{Lectura (Read Up):}\quad      & s \text{ puede observar } o \iff i(o) \ge i(s).
\end{align*}
Mnemotecnia: \emph{read up, write down}; la información ``baja'' de nivel. Intuición: ``uno es
tan confiable como las fuentes que usa'' — leo de niveles iguales o más confiables y escribo a
niveles iguales o menos confiables, para no contaminar datos de mayor integridad.

\paragraph{Notas y conexiones.} Aquí los niveles van $i(\text{Empleado}) \le
i(\text{Gerente1}), i(\text{Gerente2}) \le i(\text{CEO})$, con Gerente1 y Gerente2
\emph{incomparables} entre sí (es un orden parcial, no total: no hay camino entre ellos $\Rightarrow$
ni lectura ni escritura entre las dos ramas). El CEO es el nivel de \emph{mayor integridad/
confianza}; el Empleado el menor. Invertir las flechas cambia el objetivo de \textbf{integridad}
a \textbf{confidencialidad}: es la dualidad Biba$\leftrightarrow$Bell--LaPadula.

\paragraph{Resolución.}
\textbf{a.i) Política.} Es una política de \textbf{integridad}: el modelo \textbf{Biba}. La
etiqueta $i(\cdot)$ es nivel de integridad (confianza) y las flechas son la relación de
\textbf{dominancia} ($i(\text{Gerente1})\le i(\text{CEO})$, etc.).

\medskip
\textbf{a.ii) Propiedades y qué se puede escribir.} Las dos reglas principales de Biba
(estricto) son:
\begin{itemize}[leftmargin=1.4em,itemsep=2pt]
  \item \textbf{No Write Up (Write Down):} un sujeto $s$ puede escribir en un objeto $o$ solo
        si $i(s) \ge i(o)$. Impide que información de baja confianza contamine objetos de alta
        integridad.
  \item \textbf{No Read Down (Read Up):} un sujeto $s$ puede leer un objeto $o$ solo si
        $i(o) \ge i(s)$. Impide que un sujeto de alta integridad se ``ensucie'' leyendo datos
        poco confiables.
\end{itemize}
En base al modo de acceso \texttt{w}, cada usuario \textbf{escribe hacia abajo o a su mismo
nivel} (nunca hacia arriba). Concretamente, siguiendo el retículo:
\begin{itemize}[leftmargin=1.4em,itemsep=1pt]
  \item El \textbf{CEO} (máxima integridad) puede escribir objetos de CEO, Gerente1, Gerente2 y
        Empleado (todos los de su nivel o inferior).
  \item Un \textbf{Gerente} puede escribir objetos de su propio nivel y del Empleado, pero
        \textbf{no} del CEO (no write up) ni del otro gerente (incomparables).
  \item El \textbf{Empleado} (mínima integridad) solo puede escribir objetos de nivel Empleado;
        no puede escribir nada de los gerentes ni del CEO.
\end{itemize}
Las flechas marcan la dirección permitida del \emph{flujo} (write up está prohibido), de modo
que la información de menor confianza nunca sube a contaminar la de mayor confianza.

\medskip
\textbf{b) Invertir las flechas.} Si se invierte la relación, el modelo pasa a hacer cumplir
\textbf{confidencialidad}: es \textbf{Bell--LaPadula}, el dual de Biba. Allí la etiqueta es un
nivel de \emph{secreto/clasificación} y las reglas se invierten: \textbf{No Read Up} (seguridad
simple: leer hacia abajo, $L(s)\ge L(o)$) y \textbf{No Write Down} (propiedad $\star$: escribir
hacia arriba, $L(o)\ge L(s)$). Comparación: Biba protege la \emph{integridad} (la info ``baja'',
write-down/read-up, para no corromper lo confiable); Bell--LaPadula protege la
\emph{confidencialidad} (la info ``sube'', write-up/read-down, para no filtrar lo secreto). Son
estructuralmente inversos: ambos son control de \emph{flujo} sobre el mismo tipo de retículo,
pero con las direcciones de lectura/escritura intercambiadas.

% =====================================================================
\subsection{Ejercicio 4 — Entropía en audio WAV y esteganografía}
% =====================================================================

\begin{consigna}
En el formato de audio WAV de 32 bits, el muestreo del audio se almacena en un arreglo de
\texttt{floats} (es decir, cada muestra ocupa 4 bytes).\footnote{Los números de punto flotante
reservan valores para representar NaNs y otros valores especiales (como infinito). Ignorar esto
y asumir que todos los valores son posibles.}
\begin{enumerate}[label=\alph*)]
  \item Dado una muestra (un elemento de 4 bytes) del audio, calcular la entropía máxima en bits
        del mismo. Dar un ejemplo de un audio donde la entropía de sus muestras sea mínima.
  \item Dados archivos con 512000 muestras, se desea usar este formato para implementar un
        esquema de esteganografía, y ocultar información. La información se oculta cada 100
        muestras, es decir en el \texttt{float} número 100, 200, etc., permitiendo ocultar un
        máximo de 20480 bytes ($5120 \times 4$). La información es texto plano formado por el
        espacio en blanco y por mayúsculas y minúsculas del alfabeto inglés (exclusivamente). Si
        el secreto es menor a 20480 bytes, se lo paddea con espacios en blanco.
        Demostrar mediante cálculos de entropía, qué característica distintiva tendrían los bytes
        de audios esteganografiados con respecto a los no alterados.
  \item Postular de manera práctica cómo podrían usar este método para identificar cuáles bytes
        son alterados.
\end{enumerate}
\end{consigna}

\paragraph{Temas.} Clase 9 — Flujo de Información y Malware (\S Entropía de Shannon: máximo y
mínimo; canales encubiertos / esteganografía); Clase 8 — Principios de Diseño y Vulnerabilidades
(\S Fórmulas de entropía).

\paragraph{Qué necesitás saber.} Entropía de Shannon $H(X) = -\sum_i p(x_i)\,\log_2 p(x_i)$,
medida en bits. \textbf{Máxima} con distribución \emph{uniforme} sobre $n$ valores:
$H_{\max} = \log_2 n$. \textbf{Mínima} ($=0$) cuando un valor tiene $p=1$ (certeza: evento
único). Esteganografía $=$ ocultar información dentro de otro medio: es un \textbf{canal
encubierto} de almacenamiento. La clave del ejercicio: el alfabeto del secreto (espacio + 26
mayúsculas + 26 minúsculas $= 53$ símbolos) es \emph{mucho} más chico que el de un byte/float
arbitrario, así que las muestras alteradas tienen \textbf{mucha menos entropía}.

\paragraph{Notas y conexiones.} Reducir la entropía de un conjunto de bytes respecto del audio
original es exactamente la firma de un \textbf{flujo de información} oculto (Clase 9): el medio
``audio'' transporta información para la que no fue diseñado (canal encubierto), y medir su
entropía permite \emph{detectar} la anomalía. Es el mismo principio que el caso del disipador
que filtra bits de la clave: una vía no diseñada que baja la entropía.

\paragraph{Resolución.}
\textbf{a) Entropía máxima de una muestra.} Una muestra son $4$ bytes $= 32$ bits, es decir
$n = 2^{32}$ valores posibles. La entropía es máxima con todos equiprobables
($p = 1/2^{32}$):
\begin{equation*}
  H_{\max} = \log_2 2^{32} = \boxed{32 \text{ bits}}.
\end{equation*}
\textbf{Audio de entropía mínima:} un audio en el que todas las muestras valen \emph{lo mismo}
(p.\ ej.\ silencio total, todas las muestras $=0$, o un tono constante). Entonces un único valor
tiene $p=1$ y
\begin{equation*}
  H = -1 \cdot \log_2 1 = 0 \text{ bits} \quad (\text{certeza, no hay incertidumbre}).
\end{equation*}

\medskip
\textbf{b) Característica distintiva de los bytes esteganografiados.} En las muestras alteradas
(los floats $100, 200, \dots$) ya no se guarda un valor de audio arbitrario sino un carácter del
secreto. El alfabeto es: espacio en blanco $+$ $26$ mayúsculas $+$ $26$ minúsculas $= 53$
símbolos posibles. Si los suponemos equiprobables, cada \emph{símbolo} oculto aporta a lo sumo
\begin{equation*}
  H_{\text{secreto}} = \log_2 53 \approx 5{,}73 \text{ bits},
\end{equation*}
frente a los $32$ bits de una muestra de audio sin alterar. Más aún, el padding con espacios en
blanco \textbf{baja todavía más} la entropía efectiva (muchas muestras repiten el mismo carácter
``espacio'', acercándose a $p=1$ y por lo tanto a $H \to 0$). Conclusión:
\begin{equation*}
  H_{\text{muestras alteradas}} \;\le\; \log_2 53 \approx 5{,}73 \text{ bits}
  \;\;\ll\;\; H_{\text{muestras originales}} \le 32 \text{ bits}.
\end{equation*}
\textbf{La característica distintiva es una entropía mucho menor} en los bytes que cargan el
secreto: solo toman $53$ valores (y aún menos por el padding), mientras que el audio real ocupa
un rango amplio del espacio de $2^{32}$ valores. Esa caída de entropía \emph{es} la huella del
flujo de información oculto (canal encubierto).

\medskip
\textbf{c) Identificar prácticamente los bytes alterados.} Recorrer el arreglo de muestras y
\textbf{medir la entropía (o la diversidad de valores) por bloques}, comparando las posiciones
periódicas con el resto:
\begin{itemize}[leftmargin=1.4em,itemsep=2pt]
  \item Calcular la distribución/entropía de las muestras en las posiciones $100,200,300,\dots$
        (cada $100$) y compararla con la del resto del audio. Las posiciones del secreto
        mostrarán entropía marcadamente menor y un conjunto de valores restringido a $\le 53$
        patrones (coincidentes con códigos de letras/espacio).
  \item Equivalentemente, buscar las muestras cuyos $4$ bytes decodifican a caracteres ASCII
        imprimibles del subconjunto $\{$espacio, A--Z, a--z$\}$: en un audio genuino los floats
        rara vez caen exactamente en esos $53$ valores, y mucho menos de forma \emph{periódica}
        cada $100$ muestras.
  \item La \textbf{periodicidad} (cada $100$ muestras) más la \textbf{baja entropía} en esas
        posiciones delatan el esquema; un atacante/analista que sospeche del método mide la
        entropía de las posiciones candidatas y confirma la presencia del mensaje oculto.
\end{itemize}

% =====================================================================
\subsection{Ejercicio 5 — Opción correcta (DMZ, Von Neumann/inyección, Muralla China)}
% =====================================================================

\begin{consigna}
Elegir la opción correcta y justificar los falsos en una oración.

\textbf{1- En una DMZ se localizan}
\begin{enumerate}[label=(\alph*)]
  \item Los servidores de bases de datos con información sensible que se desea resguardar.
  \item Las claves de acceso para encriptar y desencriptar información sensible aprovechando que
        es una zona militar.
  \item Los servidores Web y de Email que requieren recibir y enviar información e interactuar
        con la red externa.
\end{enumerate}

\textbf{2- El sistema operativo iOS rompe con el modelo de Von Neumann donde los datos y el
código de ejecución se encuentran en el mismo espacio de memoria al ejecutar un programa. ¿Qué
tipo de vulnerabilidad estaría involucrada?}
\begin{enumerate}[label=(\alph*)]
  \item Soluciona un típico problema de inyección de código que podría estar en los datos.
  \item Permite resolver problemas de buffer overflow.
  \item Evita problemas de CSRF.
\end{enumerate}

\textbf{3- El modelo de Seguridad de Muralla China es adecuado para}
\begin{enumerate}[label=(\alph*)]
  \item Compartimentalizar las incumbencias para evitar situaciones de perdida de integridad en
        la información de los COIs.
  \item Compartimentalizar las incumbencias para evitar situaciones de perdida de
        confidencialidad en la información de los CDs.
  \item Compartimentalizar las incumbencias para evitar situaciones de conflicto de interes.
  \item Compartimentalizar las incumbencias para evitar situaciones de confidencialidad en los
        Company Datasets.
\end{enumerate}
\end{consigna}

\paragraph{Temas.} Clase 10 — Seguridad en la Empresa (\S Defense-in-depth, DMZ);
Clase 8 — Principios de Diseño y Vulnerabilidades (\S inyección de código; Von Neumann /
datos $=$ código; CSRF); Clase 6 — Políticas y Control de Acceso (\S Chinese Wall / Muralla
China).

\paragraph{Qué necesitás saber.}
\begin{itemize}[leftmargin=1.4em,itemsep=2pt]
  \item \textbf{DMZ:} zona \emph{expuesta} entre el firewall externo y el interno; aloja los
        servicios que \emph{deben} hablar con la red externa (Web, Email, DNS público), no los
        datos sensibles (esos van en la capa más interna, ``Datos'').
  \item \textbf{Von Neumann / inyección:} cuando datos y código comparten el mismo espacio de
        memoria, datos provistos por el usuario pueden terminar \emph{ejecutándose} como código
        $\Rightarrow$ inyección de código (la raíz del $\sim$90\% de las vulnerabilidades).
        Separar el canal de datos del de control elimina esa clase de ataque.
  \item \textbf{Muralla China (Chinese Wall):} modelo de \textbf{conflicto de interés}; objetos
        agrupados en clases de conflicto (COI), y dentro datasets de compañías (CD). Quien
        accedió a una compañía de un COI no puede acceder a una competidora del mismo COI.
\end{itemize}

\paragraph{Notas y conexiones.} El punto 2 es el principio de \emph{mecanismos exclusivos /
separación de canales} (Clase 8): el problema de fondo de toda inyección (SQLi, XSS, command
injection y hasta el prompt injection en LLMs) es mezclar el canal de \emph{datos} con el de
\emph{control}. La DMZ (punto 1) es la capa de perímetro/red de Defense-in-depth (Clase 10).

\paragraph{Resolución.}
\textbf{1) Correcta: (c).} En la DMZ se ubican los servidores Web y de Email, que por su función
deben interactuar con la red externa, manteniéndolos aislados de la red interna. \emph{Falsos:}
(a) los servidores de bases de datos con información sensible van en la capa \textbf{interna}
(``Datos''), no en la zona expuesta; (b) las claves sensibles jamás se guardan en la zona más
expuesta de la red (``DMZ'' es ``zona desmilitarizada'', un nombre, no implica resguardo
reforzado; al contrario, es la zona más expuesta).

\medskip
\textbf{2) Correcta: (a).} Separar el código del espacio de datos (no ejecutar lo que está en la
región de datos, \emph{p.\ ej.}\ con W$\oplus$X / DEP) soluciona problemas de \textbf{inyección
de código} que podría venir en los datos. \emph{Falsos:} (b) no ``resuelve'' los buffer overflow
en sí (el overflow puede seguir corrompiendo memoria/datos y causar crashes), solo evita que el
contenido inyectado se \emph{ejecute}; (c) CSRF es un ataque de confianza de sesión en la web,
no tiene relación con el modelo de memoria de Von Neumann.

\medskip
\textbf{3) Correcta: (c).} La Muralla China es un modelo de \textbf{conflicto de interés}: busca
evitar precisamente las situaciones de conflicto de interés (que un mismo sujeto acceda a datos
de empresas competidoras del mismo COI). \emph{Falsos:} (a) no es un modelo de integridad —no
protege integridad de los COI— sino de confidencialidad orientada a conflicto de interés; (b) y
(d) no apunta a ``perder/preservar confidencialidad de los CD'' como objetivo en sí, sino a
impedir el \emph{conflicto de interés}; aislar competidores es el medio, evitar el conflicto de
interés es el fin.
